%!TEX program = xelatex
\documentclass[12pt]{article}

% ---------- IDIOMA ----------
% --- Paquetes para XeLaTeX y codificación moderna ---
\usepackage{fontspec} 
\usepackage[spanish, es-tabla]{babel} 
% --- Cambiar "Cuadro" por "Tabla" ---
\addto\captionsspanish{\renewcommand{\tablename}{Tabla}}

% ---------- FUENTES ----------
\setmainfont{TeX Gyre Termes}

% ---------- PAQUETES MATEMÁTICOS Y FÍSICOS ----------
\usepackage{amsmath,amssymb}
\usepackage{amsfonts}
\usepackage{physics}

% ---------- GRÁFICOS Y COLORES ----------
\usepackage{xcolor}
\usepackage{graphicx}
\usepackage{tikz}
\usetikzlibrary{babel} % <- Agrega esta línea
\usepackage{tikz-3dplot}
\usetikzlibrary{shapes, arrows, positioning, calc}
\usetikzlibrary{arrows.meta} 
\usetikzlibrary{shapes.geometric} 
\usetikzlibrary{decorations.markings}

% ---------- CÓDIGO (LISTINGS) ----------
\usepackage{listings}
\usepackage{lastpage}

\definecolor{codegreen}{rgb}{0,0.6,0}
\definecolor{codegray}{rgb}{0.95,0.95,0.95}
\definecolor{codepurple}{rgb}{0.58,0,0.82}
\definecolor{backcolour}{rgb}{0.96,0.96,0.96}

\lstdefinestyle{pythonstyle}{
    backgroundcolor=\color{backcolour},   
    commentstyle=\color{codegreen},
    keywordstyle=\color{magenta},
    numberstyle=\tiny\color{gray},
    stringstyle=\color{codepurple},
    basicstyle=\ttfamily\small,          
    breakatwhitespace=false,         
    breaklines=true,                     
    captionpos=b,                    
    keepspaces=true,                 
    numbers=left,                        
    numbersep=5pt,                  
    showspaces=false,                
    showstringspaces=false,
    showtabs=false,                  
    tabsize=4,
    language=Python                      
}
\lstset{style=pythonstyle}

% ---------- DISEÑO Y ESTILO ----------
\usepackage{geometry}
\geometry{letterpaper, margin=2cm, headheight=15pt}
\usepackage{fancyhdr} 
\usepackage{lastpage} 
\usepackage{wrapfig}
\usepackage{multicol}
\usepackage{enumitem}
\usepackage{framed}
\usepackage{etoolbox}

% ---------- TCOLORBOX ----------
\usepackage[most,skins,theorems]{tcolorbox}

% --- Caja para Definiciones Formales ---
\newtcolorbox{definicion}[1][]{
  enhanced,
  breakable, 
  colback=blue!5!white,
  colframe=blue!75!black,
  fonttitle=\bfseries,
  title={Definición: #1},
  arc=2mm,
  boxrule=0.8pt,
  drop shadow=black!10!white
}

% --- Caja para Tareas ---
\newtcolorbox{problemas}[1][]{
  enhanced,
  breakable, 
  colback=violet!6!white,       
  colframe=violet!45!white,     
  coltitle=violet!40!black,     
  fonttitle=\bfseries,
  title={Problemas: #1},
  arc=2.5mm,                    
  boxrule=0.8pt,
  drop shadow=black!8!white     
}

% --- Caja Libre Verde Menta ---
\newtcolorbox{cajaflex}[1][]{
  enhanced,
  breakable, 
  colback=teal!5!white,       
  colframe=teal!50!black,     
  coltitle=teal!10!white,     
  fonttitle=\bfseries,
  code={\ifblank{#1}{}{\tcbset{title={#1}}}},
  arc=2.5mm,                  
  boxrule=0.8pt,
  drop shadow=black!8!white   
}

% ---------- TABLAS Y ELEMENTOS EXTRA ----------
\usepackage{booktabs}   
\usepackage{pifont}     
\usepackage{array}      
\usepackage{caption}    
\newcommand{\lapiz}{\ding{46}\ }

% ---------- ENLACES Y PDF ----------
\usepackage{hyperref}

\definecolor{myblue}{RGB}{20, 20, 120}
\definecolor{myred}{RGB}{220, 30, 30}
\definecolor{mygreen}{RGB}{0, 130, 60}

% Personalización de nombres para \autoref
\def\tableautorefname{Tabla}
\def\figureautorefname{Figura}
\def\equationautorefname{Ecuación}
\def\sectionautorefname{Sección}

% ---------- CONFIGURACIÓN DE ENCABEZADOS (ITM) ----------
\newcommand{\logo}{\includegraphics[height=3cm]{escudoitm}}
\newcommand{\escola}{Facultad de Ciencias Exactas y Aplicadas}
\newcommand{\disc}{Física Computacional 1}
\newcommand{\auth}{Robinson Longas Bedoya}
\newcommand{\emaill}{\url{robinsonlongas@itm.edu.co}}
\newcommand{\email}{\url{robinsonlongas3586@correo.itm.edu.co}}
\newcommand{\cabec}{Física Computacional 1}

\hypersetup{
    colorlinks=true,
    linkcolor=blue,
    urlcolor=blue,
    citecolor=blue,
    pdftitle={\cabec\ -- \auth},
    pdfauthor={\auth}
}

% Configuración de estilos de página
\pagestyle{fancy}
\fancyhf{} 

\fancypagestyle{firstpage}{
    \fancyhead[L]{\logo}
    \fancyhead[R]{
        \textsf{\escola \\} 
        \textsf{\disc \\}
        \textsf{\auth \\}
        \textsf{\email \\} 
        \textsf{\emaill}     
    }
    \fancyfoot[R]{\scriptsize Página \thepage~de \pageref{LastPage}}
    \renewcommand{\headrulewidth}{0pt} 
    \renewcommand{\footrulewidth}{0.4pt} 
}

\fancyhead[L]{\textsf{\scriptsize \cabec\ -- \auth\ -- \email\ -- \emaill}}
\fancyfoot[R]{\scriptsize Página \thepage~de \pageref{LastPage}}
\renewcommand{\headrulewidth}{0.4pt} 
\renewcommand{\footrulewidth}{0.4pt} 

\setlength{\parskip}{0.3em}

\begin{document}
		\thispagestyle{firstpage}
		\vspace*{1.5cm}
	
	\begin{center}
		\rule{\textwidth}{0.4pt}
	\end{center}
	
	\begin{itemize}[leftmargin=*,label=\textcolor{red}{\ensuremath{\blacktriangleright}}]
		\item \textbf{Funciones}
		\item \textbf{Módulos}
	\end{itemize}
	
	\noindent
	\textbf{\underline{Bibliografía}}
	\begin{itemize}[leftmargin=*, label={\textcolor{blue}{$\bullet$}}]
		\item \textit{Computational Physics, Capítulo 2}, Mark Newman.
		\item \textit{Introducción a la programación con Python. Capítulo 6}, Andrés Marzal, Isabel Gracia y Pedro García.
        \item \textit{Programming for computations- Python,  Capítulo 4 y sección: 5.4}, Evein Linge and Petter Langtangen. 
        \item \textit{Computational mathematics: An introduction to Numerical Analysis and Scientific Computing  with Python.
            Sección: 1.5}, Dimitrios Mitsotakis. 
    \end{itemize}
     
    

\section*{\textcolor{red}{\textit{Funciones}}}
	Hasta el momento, utilizamos las funciones  que traen definidas las 
    librerías de Python. Por ejplo,  \lstinline|sqrt|, \lstinline|cos|, 
    \lstinline|arcos|, por nombrar unas cuantas. Sin embargo, existen muchas 
    situaciones en Física Computacional, en las que necesitamos definir nuestras
    propias funciones.
    
    La estructura básica para escribir funciones en Python es:
    %
	\begin{lstlisting}
def la_funcion(parametro1,parametro2=valor2,...):
	''' Ayuda de la función '''
	declaraciones
	return valor_funcion
	\end{lstlisting}
	%
	El texto que se escriba en la '''Ayuda de la función''' sirve para indicarnos luego 
	de que se trata la función y queda definido como el atributo del objeto: 
	\lstinline|la_funcion.__doc__| o \lstinline|la_funcion?| en Ipython.
	El argumento opcional valor2 es el valor que adquiere el parámetro si la 
	función es llamada sin especificar el parametro2. 

    Por ejemplo, $f(x) = 8x^7 - x^3 + x^2 - 1$,
    %
\begin{lstlisting}[language=Python]
In [1]: def f(x):
   ...:     y = 8*x**7 - x**3 + x**2 -1
   ...:     return y
   ...: 

In [2]: f(0)
Out[2]: -1

In [3]: f(-1)
Out[3]: -7

In [4]: def f(x):
   ...:     return 8*x**7 - x**3 + x**2 -1
   ...: 

In [5]: f(-1)
Out[5]: -7

In [6]: f(0)
Out[6]: -1
\end{lstlisting}
% 
Note que se puede omitir la definición de la variable $y$ poniendo la operación después del \lstinline|return|.
En ocasiones, queremos entregar un valor por defecto a alguna de las variables. Por ejemplo,
\begin{lstlisting}[language=Python]
In [1]: def f(x):
   ...:     return 8*x**7 - x**3 + x**2 -1
   ...: 

In [2]: f(-1)
Out[2]: -7

In [3]: f(0)
Out[3]: -1
In [4]: def maximo(x,y=0):
   ...:     if x >=y:
   ...:         return x
   ...:     else:
   ...:         return y
   ...: 

In [5]: maximo(1)
Out[5]: 1

In [6]: maximo(2)
Out[6]: 2

In [7]: maximo(2),5
Out[7]: (2, 5)

In [8]: maximo(2,5)
Out[8]: 5

In [9]: maximo(7,5)
Out[9]: 7

In [10]: maximo(0,5)
Out[10]: 5

In [11]: maximo(0)
Out[11]: 0

In [12]: maximo(-1)
Out[12]: 0

In [13]: a=50; b=100;

In [14]: maximo(a,b)
Out[14]: 100
\end{lstlisting}
% 
    
Como último ejemplo, queremos calcular el factorial de un número, 
    $$ n! = 1 \times 2 \times 3 \times ....( n - 2) \times (n - 1) \times n = \prod_{k=1}^{n} k \,.$$
    
    Es decir, $3! = 1\times 2\times 3 = 6$.
%
\begin{lstlisting}[language=Python, captionpos=t, title=\texttt{factorial.py}]
""" Este es un ejemplo del uso de funciones """

def factorial(n):
    ''' Calcula el factorial de un número entero'''
    f = 1
    for i in range(1,n+1):
        f = i * f
    return f

### acá comenzamos el programa

a = int(input("ingrese el número entero del que quiere conocer el factrorial\n"))

if a>=0:
    elfatorial = factorial(a)
    print(f"El factorial del número que entraste es: {elfatorial} ")
else:
    print("No es la función Gamma, el número debe ser entero mayor o igual a cero.")
\end{lstlisting}
%  

El ejemplo del cálculo del factorial nos permite analizar paso a paso los componentes esenciales 
de una función en Python y la dinámica de ejecución dentro del programa:

\begin{enumerate}
    \item \textbf{Definición de la función (\lstinline|def|):} 
    La instrucción \lstinline|def factorial(n):| le indica al intérprete que estamos creando un bloque de código reutilizable llamado \lstinline|factorial|, el cual espera recibir un argumento llamado \lstinline|n| (el parámetro de entrada).

    \item \textbf{Documentación o Docstring:} 
    La cadena dentro de triples comillas \lstinline|'''Calcula el factorial de un número entero'''| constituye el \textit{docstring}. Es una práctica fundamental en software científico, ya que permite consultar la ayuda del objeto mediante \lstinline|factorial?| o \lstinline|help(factorial)| en entornos interactivos como IPython o Jupyter.

    \item \textbf{Variables locales y acumuladores:} 
    Dentro de la función se define \lstinline|f = 1|. Esta variable es de \textbf{ámbito local} (\textit{local scope}), lo que significa que solo existe mientras la función se está ejecutando. El bucle \lstinline|for i in range(1, n + 1):| actualiza el acumulador multiplicando \lstinline|f = i * f| (o equivalentemente \lstinline|f *= i|) desde $i = 1$ hasta $i = n$.

    \item \textbf{Valor de retorno (\lstinline|return|):} 
    La sentencia \lstinline|return f| finaliza la ejecución de la función y devuelve el resultado computado al punto exacto del programa donde fue invocada. Sin un \lstinline|return|, la función devolvería \lstinline|None|.

    \item \textbf{Flujo del programa principal y paso de parámetros:} 
    En la sección principal del script:
    \begin{itemize}
        \item Se captura la entrada del usuario mediante \lstinline|input()| y se convierte explicitamente a entero con \lstinline|int()|.
        \item Se evalúa mediante una estructura condicional (\lstinline|if a >= 0|) que el número sea no negativo.
        \item En la línea \lstinline|elfactorial = factorial(a)|, la variable externa \lstinline|a| pasa su valor al parámetro local \lstinline|n|. El resultado devuelto por \lstinline|return| se almacena en la nueva variable \lstinline|elfactorial|.
        \item Finalmente, se despliega el resultado mediante una 
        cadena formateada: \lstinline|f"El factorial del número que entraste es: {elfactorial}"|.
    \end{itemize}
\end{enumerate}

\begin{cajaflex}[Observación]
Note que para usar la función debo hacerlo afuera de su definición, después del comentario 
\lstinline|### acá comienza el programa|. Python permite el uso la estructura 
\lstinline|if__name__=="__main"":|. Lo veremos después cuando hablemos de módulos.
\end{cajaflex}

\subsection*{\textcolor{blue}{\textit{Ejemplo:}}}
\begin{lstlisting}[language=Python, captionpos=t, title=\texttt{tiempo\_vuelo\_v3.py}]
import numpy as np
import matplotlib.pyplot as plt


def altura(y0, v0, t):
    """Esta función entrega la altura de una bola en función del tiermpo'
    Parametros: altura inicial y0, velocidad inicial v0 y tiempo"""
    g = 9.8 # aceleración gravitacional de la tierra en m/s^2
    return  y0 + v0 * t - 1/2. * g * t**2 


#--- programa principal ---------------
print('----------------------------------')
y0 = float(input('Entre la altura de lanzamiento en metros\t'))
v0 = float(input('Entre la rapidez con la que lo lanzas en m/s\t'))
N = int(input("Entre el número de puntos para el paso del tiempo\t")) 
print('----------------------------------')
g = 9.8 # aceleración gravitacional de la tierra en m/s^2

#--------- solucion exacta ------------
T = (v0 + np.sqrt( v0**2 + 2*g*y0) )/g
#------------------------

t_i = 0.0 # tiempo inicial
t_f = 10. #tiempo final
t = np.linspace(t_i, t_f, N) # arreglo con los tiempos

# llamos a la función altura
y = altura(y0,v0,t) 

# con el while buscamos que esa altura siempre sea positiva. 
i = 0
while y[i] >=0:
    i = i +1

t_num=t[i]

# calculo de errores

err_abs = abs(t_num - T)
err_rel = (err_abs / T) * 100    

print(f"El tiempo que queda la bola en el aire: {t_num:.4f} segundos")
print(f"El tiempo exacto en el que cae la bola es : {T:.4f} segundos")
print("\n--- ANÁLISIS DE ERROR Y PRECISIÓN ---")
print(f"Error absoluto (|t_num - T|)  : {err_abs:.6f} s")
print(f"Error relativo (%)            : {err_rel:.4f} %")
print(f"La altura para ese tiempo es: {y[i]:.4e} metros")
\end{lstlisting}
%  
Note que la aceleración gravitacional se definió en el interior de la función, esto la 
convierte en una variable local. Eso significa que, para calcular la solución 
exacta \lstinline|T = (v0 + np.sqrt( v0**2 + 2*g*y0) )/g| debo volver a definir
$g$ afuera de la función. Mucho cuidado con el orden en que se entran los argumentos de la función,
\lstinline|altura(v0,y0,t)| no genera error pero devuelve un resultado incorrecto
fisicamente. 

Para ilustrar el concepto de las variables locales y globales, considere el siguiente ejemplo
\begin{lstlisting}[language=Python]
In [16]: def variables(x,y):
    ...:     temp = x
    ...:     x = y
    ...:     y= temp
    ...:     print('Estas son variables locales: x =',x, 'y=',y)
    ...:     return x,y

In [17]: x = 2; y= 3

In [18]: variables(x,y);
Estas son variables locales: x = 3 y= 2

In [19]: print("Estas son las variables globales: x=",x,'y=',y)
Estas son las variables globales: x= 2 y= 3
\end{lstlisting}
% 

\subsection*{\textcolor{blue}{\textit{La función Lambda}}}
En ocasiones simplemente queremos definir una fiunción corta, como por ejemplo
$f(x) = x^3 - x$ o $g(x)= x^2$.
\begin{lstlisting}
In [1]: g = lambda x: x**2

In [2]: g(-1)
Out[2]: 1

In [3]: g(1)
Out[3]: 1

In [4]: g(45)
Out[4]: 2025

In [5]: def g(x): # equivalente al uso de lambda
   ...:     return x**2
   ...:     
\end{lstlisting}

\section*{\textcolor{red}{\textit{Módulos en Python}}}
Un módulo es un archivo que contiene un conjunto de funciones y su estructura 
	está dada por la siguiente sintaxis:
	%
	\begin{lstlisting}
def la_funcion1(parametro1,parametro2=valor2,...):
	''' Ayuda de la función '''
	declaraciones
	return valor_funcion

def la_funcion2(parametro1,parametro2=valor2,...):
	''' Ayuda de la función '''
	declaraciones
	return valor_funcion

def la_funcion3(parametro1,parametro2=valor2,...):
	''' Ayuda de la función '''
	declaraciones
	return valor_funcion
......

if __name__ == "__main__":
	Código de prueba o test del módulo
	\end{lstlisting}
	%
	En Python viene una variable definida internamente \lstinline|__name__|
	Esta variable toma el valor del programa cuando se carga como módulo o el valor
	cuando se corre como programa. 
    
    Veamos por ejemplo como se vería la función 
    para el cálculol del factorial,
%
\begin{lstlisting}[language=Python, captionpos=t, title=\texttt{factorial\_v2.py}]
""" Este es un ejemplo del uso de funciones como módulo para otros programas
Devuelve el factorial de n entero """

def factorial(n):
    ''' Calcula el factorial de un número entero n! = n*(n-1)*(n-2)...'''
    f = 1.
    for i in range(1,n+1):
        f = i * f
    return f


#----- programa principal -------
if __name__ == "__main__":
    a = int(input("ingrese el número entero del que quiere conocer el factrorial\n"))

    if a>=0:
        elfatorial = factorial(a)
        print(f"El factorial del número que entraste es: {elfatorial} ")
    else:
        print("No es la función Gamma, el número debe ser entero mayor o igual a cero.")
\end{lstlisting}
%  

\subsection*{\textcolor{blue}{\textit{El Punto de Entrada: \texttt{if \_\_name\_\_ == "\_\_main\_\_":}}}}

En nuestros programas es fundamental que los de código sean 
modulares, es decir, que puedan funcionar tanto como programas 
independientes que se ejecutan directamente, como librerías de funciones 
importables desde otros archivos. 

Para lograr este comportamiento, Python utiliza una variable especial interna llamada \lstinline|__name__|.

\subsubsection*{1. ¿Qué es \texttt{\_\_name\_\_} y cómo funciona?}

Cada vez que el intérprete de Python ejecuta un archivo de código (\textsf{.py}), asigna automáticamente un valor a la variable oculta \lstinline|__name__| según el contexto:

\begin{itemize}
    \item \textbf{Si el archivo se ejecuta directamente} (por ejemplo, corriendo \texttt{python factorial.py} en la terminal):
    Python le asigna a \lstinline|__name__| la cadena exacta \lstinline|"__main__"|.
    
    \item \textbf{Si el archivo es importado como módulo} dentro de otro programa (por ejemplo, haciendo \lstinline|import factorial| desde otro script):
    Python le asigna a \lstinline|__name__| el nombre real del archivo (en este caso, \lstinline|"factorial"|).
\end{itemize}

\subsubsection*{2. El propósito del condicional}

Al escribir la estructura:

\begin{lstlisting}
if __name__ == "__main__":
    # Codigo del programa principal
    a = int(input("Ingrese un numero: "))
    print(factorial(a))
\end{lstlisting}

Le estamos diciendo a Python: 
\begin{quote}
\textit{«Ejecuta este bloque de código \textbf{únicamente} si el archivo está 
siendo ejecutado de forma directa. Si el archivo está siendo importado 
por otro script, ignora este bloque y limita tu trabajo a ofrecer 
las funciones definidades arriba.»}
\end{quote}

Veamos un ejemplo utilizando la función factorial para calcular los coeficientes 
binomiales 
$$ \binom{n}{k} = \frac{n!}{k!(n - k)!} \quad \text{para } 0 \le k \le n \,,$$
y el triángulo de Pascal.
%
\begin{lstlisting}[language=Python, captionpos=t, title=\texttt{triangulo\_pascal.py}]
""" Lee la función factorial, calcula la combinatoria y muestra el triángulo 
de Pascal""" 

import factorial_v2  as fa # importamos nuestra función

def binomial(n, k):
    """
    Calcula el coeficiente binomial C(n, k) = n! / (k! * (n - k)!)
    usando la función factorial previamente definida.
    Retorna un número entero.
    """
    if k < 0 or k > n:
        return 0
    if k == 0 or k == n:
        return 1
    
    # Usamos la función factorial que llamamos 
    num = fa.factorial(n)
    den = fa.factorial(k) * fa.factorial(n - k)
    return int(num/den)

def imprimir_pascal(filas):
    """Imprime las primeras filas del Triángulo de Pascal."""
    for n in range(filas):
        linea = []
        for k in range(n + 1):
            # Llamamos a la función de coeficientes binomial(n, k)
            coeficiente = binomial(n, k)
            linea.append(str(coeficiente))
        
        # Unir los números con espacios sin centrar
        #print(" ".join(linea))
        # Centra el texto en un ancho fijo de 40 caracteres
        print(" ".join(linea).center(40))

if __name__ == "__main__":

    # Pedir las filas que desee del triángulo de Pasca
    n = int(input("¿Cuántas filas desea para el triángulo?\t"))
    imprimir_pascal(n)
\end{lstlisting}
%  

\begin{enumerate}
    \item \textbf{Importación de Módulos Propios y Alias (\lstinline|import ... as|):}
    \begin{itemize}
        \item La instrucción \lstinline|import factorial_v2 as fa| busca el archivo \texttt{factorial\_v2.py} en el mismo directorio de trabajo e importa todas sus funciones.
        \item El uso del alias \lstinline|as fa| nos permite invocar la función con una sintaxis más corta y limpia: \lstinline|fa.factorial(n)| en lugar de \lstinline|factorial_v2.factorial(n)|.
    \end{itemize}

    \item \textbf{Definición de \lstinline|binomial(n, k)|:}
    \begin{itemize}
        \item Se incluyen optimizaciones con casos base: si $k < 0$ o $k > n$, el coeficiente es $0$; si $k = 0$ o $k = n$, devuelve directamente $1$ sin necesidad de calcular factoriales ($\binom{n}{0} = \binom{n}{n} = 1$).
        \item \textbf{Nota de precisión:} La línea \lstinline|return int(num / den)| o idealmente la división entera \lstinline|return num // den| garantiza que el resultado sea de tipo entero (\lstinline|int|), evitando decimales innecesarios en la impresión (e.g., $6$ en lugar de $6.0$).
    \end{itemize}

    \item \textbf{Construcción de la Cadena :}
    \begin{itemize}
        \item \lstinline|linea.append(str(coeficiente))|: Convierte cada número que se tiene en la variable \newline \lstinline|coeficiente|
        en cadena para almacenarlo en la lista temporal \lstinline|linea|.
        \item \lstinline|" ".join(linea)|: Concatena los elementos de la lista intercalando un espacio entre ellos (ejemplo: \texttt{['1', '3', '3', '1']} $\rightarrow$ \texttt{"1 3 3 1"}).
        \item \lstinline|.center(40)|: Es un método de cadenas de Python que añade espacios en blanco de forma equitativa a la izquierda y derecha para centrar la fila dentro de un bloque fijo de $40$ caracteres de ancho, generando la geometría piramidal clásica del triángulo.
    \end{itemize}
\end{enumerate}

\subsection*{\textcolor{blue}{\textit{Ejemplo: secuencia de Fibonacci}}}

Veamos un ejemplo de un módulo que me genera los números de Fibonacci. 

	\noindent \textcolor{blue}{\textit{La siguiente discusión fue tomada del libro: 
	Introduction to real analysis, Robert G. Bartle et al}}

	La secuencia de Fibonacci es una sucesión infinita de números 
	naturales que comienza con 1, y cada número subsiguiente es
	la suma de los dos anteriores. 
	La secuencia se define de la siguiente manera, 
	$F:=(f_n)_{n \in \mathbb{N}}$
	\begin{align}\label{eq:sucesion_fibo}
		f_1:=1,\hspace{1cm}f_2:=1,\hspace{1cm}f_{n+1}:=f_{n-1}+f_{n},\;(n\geq 2)\,.
	\end{align}
	Note que los primeros números de Fibonacci son: 1, 1, 2, 3, 5, 8, 13...
	Dentro de muchas propiedades de la sucesión de Fibonacci, 
	se encuentra que el cociente de dos números consecutivos tiende a la razón
	áurea $\phi=\frac{1+\sqrt{5}}{2}\approx 1.6180339887$ (para más detalles 
	ver el siguiente video: \href{https://www.youtube.com/watch?v=yDyMSliKsxI}{La sucesión de fibonacci y la razón aúrea})  
	Acá te dejo otros dos videos interesantes sobre la sucesión de Fibonacci:
	\href{https://www.youtube.com/watch?v=N0gr6nK6lgU}{Lo que nadie sabe de la sucesión 
	de Fibonacci} y \href{http://youtube.com/watch?v=ctRIK56jV3k}{Las plantas y fibonacci}. 
	También te dejo un \href{https://www.bbc.com/mundo/noticias-46926506}{artículo de divulgación}.

	Vamos entonces a probar que $\lim_{n\rightarrow \infty} \dfrac{f_{n+1}}{f_n} = \phi$.
	Para ello partimos de la ecuación \eqref{eq:sucesion_fibo} 
	y dividimos ambos lados por $f_n$ para obtener
	\begin{align}
		\dfrac{f_{n+1}}{f_n} = \dfrac{f_{n-1}}{f_n}+1\,.
	\end{align}
	Si definimos $x_n=\dfrac{f_{n+1}}{f_n}$, entonces la ecuación anterior se 
	puede reescribir como
	\begin{align}
		x_n = \dfrac{1}{x_{n-1}}+1\,.
	\end{align}		
	Asumiendo que la sucesión $(x_n)_{n\in \mathbb{N}}$ converge a  $L$, 
	se tiene que 
	\begin{align}
		\lim_{n\rightarrow \infty} x_n = \lim_{n\rightarrow \infty} \left[ 
			\dfrac{1}{ x_n}+1 \right] \Rightarrow  L = \frac{1}{L} + 1 \implies 
			L^2-L-1=0 \,.
	\end{align}	
	La ecuación anterior se resuelve de la manera usual y se encuentran dos raíces.
	La raíz positiva coincide con la razón aúrea $L=\dfrac{1+\sqrt{5}}{2}\equiv\phi$.

A continuación, mostraremos numéricamente el resultado anterior.
\begin{lstlisting}[language=Python, captionpos=t, title=\texttt{fibonacci.py}]
""" 
fibonacci. py
Es un módulo que genera los números de Fibonacci y se verifica la famosa razón aúrea. 
"""
import math 

def Fibonacci(N):
    """
    Genera y devuelve una lista  con los primeros N números de Fibonacci.
    
    Parametros de entrada:
    N (int): Número de Fibonacci que desea generar.
    
    Devuelve:
    list: Los primeros N números de Fibonacci comenzando desde f_1 = 1 y f_2 = 1.
    """
    if N <= 0:
        print("Error: N debe ser un entero positivo.")
        return []
    elif N == 1:
        return [1]
    
    # Inicia la secuencia (lista) con los dos primeros números 
    fib_sequence = [1, 1]
    
    # Agrego los números de Fibonacci hasta que la longitud de la lista sea N, con el .append
    # Tomamos el último de la lista fib_sequence[-1] 
    # Tomamos el penúltimo de la lista fib_sequence[-2] 
    # Agregamos cada elemento a la lista fib_sequence.append(fib_sequence[-1] + fib_sequence[-2])
    # Devuelve la lista desde e primero hasta el N-1 fib_sequence[:N]
    for i in range(2, N):
        fib_sequence.append(fib_sequence[-1] + fib_sequence[-2])
        
    return fib_sequence[:N]

def convergencia(F):
    """
    Acá chequeamos la convergencia de la sucesión de Fibonacci y comparamos 
    con la razón aúrea.
    
    Paramétro de entrada: La lista F con los números de Fibonacci.
    
    Retorna:
    None: Imprime una tabla con la comparación entre la convergencia de la 
    sucesión de fibonacci y la razón aúrea.
    """
    # La rázon aúrea: (1 + sqrt(5)) / 2
    aurea = (1 + math.sqrt(5)) / 2
    
    # Se necesitan al menos dos números de Fibonacci para calcular la razón
    valid_fibs = [x for x in F if x > 0]
    
    if len(valid_fibs) < 2:
        print("Error: La lista de entrada debe contener al menos dos números de Fibonacci.")
        return

    print(f"{'n':<5} | {'f_n':<10} | {'f_n+1':<10} | {'Razón (F_n+1 / F_n)':<22} | {'Diferencia con la razón aúrea':<30}")
    print("-" * 85)

    ratios = []
    for i in range(len(valid_fibs) - 1):
        fn = valid_fibs[i]
        fn1 = valid_fibs[i+1]
        ratio = fn1 / fn
        ratios.append(ratio)
        diff = abs(ratio - aurea)
        
        print(f"{i+1:<5} | {fn:<10} | {fn1:<10} | {ratio:<22.10f} | {diff:<30.10e}")

    # Chequea si la última diferencia es cercana a 0
    final_diff = abs(ratios[-1] - aurea)
    print("\n" + "=" * 85)
    print(f"Razón aúrea (Valor real): {aurea:.10f}")
    print(f"Última razón calculada:     {ratios[-1]:.10f}")
    print(f"Discrepancia absoluta:      {final_diff:.10e}")
    
    if final_diff < 1e-4:
        print("\nVerdicto: Johannes Kepler estaba correcto! .")
    else:
        print("\nVerdicto: La secuencia es demasiado corta para mostrar una convergencia.")


# --- Programa principal ---
if __name__ == "__main__":

    print("--- Test Block del programa fibonacci.py ---\n")
    N = int(input("Ingrese el número de términos para la secuencia "))

    print(f"Generando los primeros {N} números de Fibonacci:")

    fib_list = Fibonacci(N)
    print(fib_list)
    print("\n" + "=" * 85 + "\n")
    
    print("Verificando la convergencia :")
    convergencia(fib_list)
\end{lstlisting}

%--------------- problemas --------------------

\begin{problemas}[]
\begin{enumerate}[leftmargin=*]

\item Defina una función que devuelva el valor de la
raíz cúbica de un número dado.
\item Define una función que convierta grados centígrados
a Kelvin y a Farenheit. 
\item Defina una función que convierta grados en radianes.
\item Defina una función que daods dos parámetros, $b$ y $x$,
devuelva al logartimo en base $b$ de $x$, $\log_b(x)$. 
\item \textbf{Números perfectos:}
Se dice que un número es perfecto si es igual a la suma
de todos sus divisores, excpeto él mismo. Por ejemplo:
los divisores de 28 son 1, 2, 4, 7, 14 y 28. 
Ahora, nota que 28 = 1 + 2+ 4+ 7 + 14. Entonce s28 
es un número perfecto. Diseña una función que nos
diga si un número dado es perfecto o no.

\item \subsection*{\textcolor{red}{\textit{Construcción del módulo \texttt{cinematica.py}}}}

Para resolver el siguiente problema, \textit{lea la secciónn}

\begin{enumerate}
    \item \textbf{Parte 1: Diseño del Módulo (\texttt{cinematica.py})}
    Cree un archivo llamado \texttt{cinematica.py} y definina 
    las siguientes funciones:
    \begin{itemize}
        \item \lstinline|posicion_1d(x0, v0, a, t)|: Devuelve la posición $x(t)$. 
        Debe aceptar tanto números flotantes como arreglos de \textsf{NumPy} para 
        la variable $t$.
        \item \lstinline|velocidad_1d(v0, a, t)|: Devuelve la velocidad $v(t)$.
        \item \lstinline|tiempo_vuelo_parabolico(v0, angulo_deg, y0=0.0, g=9.8)|: 
        Recibe la velocidad inicial, el ángulo en grados (convertir internamente a radianes 
        con \lstinline|np.radians|) y la altura inicial. 
        Retorna el tiempo total que el proyectil permanece en el aire.
        \item \lstinline|alcance_maximo(v0, angulo_deg, y0=0.0, g=9.8)|: Hace uso interno de la función \lstinline|tiempo_vuelo_parabolico| para retornar el alcance horizontal total $R$.
    \end{itemize}
    \textit{Requisito:} Incluya el bloque \lstinline|if __name__ == "__main__":| al final 
    de \texttt{cinematica.py} con las pruebas necesarias para validar que sus funciones 
    entregan valores físicos coherentes 
    (por ejemplo, verificar que con $v_0 = 10\text{ m/s}$, 
    $\theta = 45^\circ$, $y_0 = 0$, el alcance es aproximadamente $10.2\text{ m}$).

    \item \textbf{Parte 2: Aplicación en script principal} En un archivo separado 
    ejecute las siguientes tareas importando su módulo 
    recién creado:
    \begin{enumerate}
        \item Un cañón situado en un acantilado a $y_0 = 35\text{ m}$ 
        dispara proyectiles a $v_0 = 40\text{ m/s}$. Varíe el ángulo de 
        lanzamiento $\theta$ desde $10^\circ$ hasta $80^\circ$ en pasos de 
        $1^\circ$.
        \item Usando las funciones de su módulo, encuentre numéricamente el 
        ángulo óptimo que maximiza el alcance horizontal $R$ y el tiempo 
        de vuelo correspondiente.
        \item Imprima los resultados en pantalla con formato profesional 
        utilizando \textit{f-strings}.
        \item Haga gráficos de posición y velocidad contra el tiempo para diferentes ángulos. Específicamente:
        \begin{itemize}
            \item Grafique $x(t)$ e $y(t)$ en función del tiempo para los ángulos $\theta \in \{15^\circ, 30^\circ, 45^\circ, 60^\circ, 75^\circ\}$.
            \item Grafique la velocidad horizontal $v_x(t)$ y la velocidad vertical $v_y(t)$ en función del tiempo para los mismos ángulos, identificando el punto donde $v_y(t) = 0$ (altura máxima).
        \end{itemize} 
    \end{enumerate}
\end{enumerate}

\item \textbf{Módulo de Geometría Vectorial (\texttt{vectores.py}):}
    Diseñe un nuevo módulo llamado \texttt{vectores.py} que contenga las siguientes funciones para vectores en 2D/3D representados como listas o arreglos de NumPy:
    \begin{itemize}
        \item \lstinline|norma(v)|: Devuelve la magnitud $|\vec{v}|$.
        \item \lstinline|producto_punto(u, v)|: Devuelve el escalar $\vec{u} \cdot \vec{v}$.
        \item \lstinline|angulo_entre(u, v)|: Devuelve el ángulo en grados entre ambos vectores.
    \end{itemize}
    Incluya casos de prueba en el condicional \lstinline|if __name__ == "__main__":| para verificar que los vectores ortogonales devuelvan un ángulo de $90^\circ$.

 \item \textbf{Caída Libre con fricción lineal:}
    La velocidad de una partícula en caída libre sujeta a resistencia del aire proporcional a la velocidad está dada por la solución analítica:
    $$ v(t) = \frac{m g}{c} \left( 1 - e^{-\frac{c}{m}t} \right) $$
    donde $m$ es la masa, $g = 9.8\text{ m/s}^2$ y $c$ es el coeficiente de arrastre.
    \begin{enumerate}
        \item Cree una función \lstinline|velocidad_terminal(m, c)| que retorne la velocidad límite $v_t = \frac{mg}{c}$.
        \item Cree la función \lstinline|velocidad_friccion(t, m, c)| que evalúe la ecuación completa.
        \item Escriba el bloque \lstinline|if __name__ == "__main__":| para pedir al usuario la masa $m$ y el coeficiente $c$, e imprimir cuánto tiempo tarda la partícula en alcanzar el $99\%$ de su velocidad terminal.
    \end{enumerate}   
\end{enumerate}

\end{problemas}


\begin{cajaflex}[Fundamentos Teóricos: Cinemática en 1D y 2D]
Antes de implementar nuestro propio módulo de cinemática, es imperativo repasar 
las ecuaciones analíticas que gobiernan el movimiento de una partícula bajo aceleración constante $a$:

\paragraph{1. Movimiento en Una Dimensión (1D) con Aceleración Constante:}
Para una partícula que se mueve a lo largo de un eje con posición inicial $x_0$, velocidad inicial $v_0$ y aceleración constante $a$:
\begin{align}
    x(t) &= x_0 + v_0 t + \frac{1}{2} a t^2 \\
    v(t) &= v_0 + a t 
\end{align}

\paragraph{2. Movimiento de Proyectiles (2D) en un Campo Gravitacional:}
Cuando un proyectil se lanza desde una posición $(x_0, y_0)$ con una rapidez inicial $v_0$ y un ángulo $\theta$ respecto a la horizontal, el movimiento se desacopla en dos ejes independientes (asumiendo aceleración gravitacional $g = 9.8\text{ m/s}^2$ apuntando hacia abajo):

\begin{itemize}
    \item \textbf{Eje Horizontal ($x$):} Movimiento a velocidad constante ($a_x = 0$).
    \begin{equation}
        v_x = v_0 \cos\theta, \quad x(t) = x_0 + (v_0 \cos\theta) t
    \end{equation}
    
    \item \textbf{Eje Vertical ($y$):} Caída libre con aceleración constante ($a_y = -g$).
    \begin{equation}
        v_y(t) = v_0 \sin\theta - g t, \quad y(t) = y_0 + (v_0 \sin\theta) t - \frac{1}{2} g t^2
    \end{equation}
\end{itemize}

\paragraph{3. Magnitudes Característas del Tiro Parabólico:}
Resolviendo $y(t) = 0$ para el tiempo de vuelo $T_{\text{vuelo}}$ y evaluándolo en $x(t)$, 
obtenemos el alcance horizontal máximo $x_{max}$:
\begin{equation}
    T_{\text{vuelo}} = \frac{v_0 \sin\theta + \sqrt{(v_0 \sin\theta)^2 + 2 g y_0}}{g}
\end{equation}
\begin{equation}
    x_{max} = x(T_{\text{vuelo}}) = x_0 + (v_0 \cos\theta) T_{\text{vuelo}}
\end{equation}
\paragraph{4. Altura máxima del tiro parabólico:}
Resolviendo $v_y(t) = 0$ para el tiempo y evaluándolo en $y(t)$, 
obtenemos la altura máxima $y_{max}$.
\end{cajaflex}

\end{document}

